
\documentclass{article} % Starts an article
\usepackage{amsmath} % Imports amsmath
\begin{document}
\date{}
\title{Tom's 2024 Macro Qualifier Question}
%\centerline{\bf Homework 3}

\maketitle
\noindent
 \textbf{Question 1.} 

\medskip
\noindent \textbf{Definitions:} 
An ``intertemporal problem'' is stated in terms of  infinite sequences indexed by time $t = [0, 1, 2, \ldots, )$.
An  ``inverse problem''  asserts an  answer and invites you to state an  associated question.
 For the ``inverse problems'' in parts A, B, and C,  your job is please to
state completely the optimum problem to which that Bellman equation corresponds.  




\medskip
\noindent
 \textbf{Part A.} Where $x \in {\bf R}$, $f >0, d >0 $,   $\beta \in (0,1)$, and $g : {\bf R} \rightarrow {\bf R}$, for what intertemporal problem is the following  a Bellman equation?
$$ w(x) =  \{- f x^2 - d(g(x)  - x)^2 + \beta w(g(x)) \}  $$


\medskip
\noindent 
\textbf{ Part B.}   Where $x \in {\bf R}$, $f >0, d >0 $,    $\beta \in (0,1)$, and $w(x)$ solves the functional equation in Part A,  for what intertemporal problem is the following  a Bellman equation?
$$ v(x) = \max_{x'} \{ - f x - d(x' - x)^2 + \beta w(x') \}
 $$


\medskip

\noindent
\textbf{ Part C.} Where $x \in {\bf R}$, $f >0, d >0 $, and   $\beta \in (0,1)$, for what intertemporal problem is the following  a Bellman equation?
$$ v(x) = \max_{x'} \{ - f x - d(x' - x)^2 + \beta v(x') \}
 $$

\medskip
\noindent 
\textbf{ Part D.} Could  the Bellman equations in parts A and B provide a useful  approach to solving the Bellman equation in Part C?

\end{document}



